\documentclass{exam}
\usepackage{../../mypackages}
\usepackage{../../macros}

% Pour le tableau
\usepackage{array}
\usepackage{longtable}
\usepackage{multirow}
\usepackage[table]{xcolor}

\title{Critères de réussite}
\author{N. Bancel}
\date{Février 2026}

% Mise en forme des solutions en bleu
\SolutionEmphasis{\color{blue}}
\renewcommand{\solutiontitle}{\noindent}

\definecolor{lavande}{RGB}{180,190,255}


\tcbuselibrary{theorems,skins,hooks}

\definecolor{myindbg}{HTML}{DFF5EC}
\definecolor{myindfr}{HTML}{1F6F4A}

\newtcolorbox{Indication}{
  enhanced,
  colback=myindbg,
  colframe=myindfr,
  arc=6mm,
  rounded corners,
  boxrule=0.6mm,
  left=3mm,
  right=3mm,
  top=2mm,
  bottom=2mm
}

% Colonne p{} alignée à gauche et qui wrappe
\newcolumntype{L}[1]{>{\raggedright\arraybackslash}p{#1}}

\begin{document}

\textbf{Collège Lycée Suger}
\hfill
\textbf{Mathématiques} \\

\textbf{Année 2025-2026}
\hfill
\textbf{Classe de 6ème} \par

{\let\newpage\relax\maketitle}

\section{Général - Geométrie}

\begin{Indication}

\begin{compactitem}
  \item Tous les dessins sont tracés au crayon à papier. Sauf indiqué, les tracés ne sont jamais faits à main levée, mais avec règle, équerre, compas, rapporteur etc.
  \item Une figure se fait en trait plein, bien appuyé
  \item Les étapes qui permettent de construire la figure sont tracés en pointillés, ou avec un trait plein moins appuyé que la figure. La figure doit bien ressortir.
  \item Légendes et codages importants
  \begin{compactitem}
    \item Les points importants sont nommés (milieu, centre, sommet, etc.)
    \item Les angles droits sont codés avec un petit carré
    \item Les segments égaux sont codés avec des petits traits (un trait pour les segments égaux entre eux, deux traits pour un autre groupe de segments égaux, etc.)
    \item Les angles égaux sont codés avec des arcs, et un ou plusieurs traits dessus quand ils sont de même mesure
\end{compactitem}
\end{compactitem}

\end{Indication}


\section{Chapitre 10 - Symétrie axiale}

\begin{Indication}

\begin{compactitem}
  \item Quand on trace le symétrique d'un objet mathématique (point, droite, cercle etc), on lui donne le même nom avec une apostrophe (exemple : le symétrique de $A$ est noté $A'$, le symétrique d'une droite $d$ est $(d')$).
  \item Un axe de symétrie est une droite, pas un segment
  \item Quand on trace le symétrique d'un point par rapport à un axe $(d)$
\begin{compactitem}

    \item On est très précis sur la perpendicularité (faite à l'équerre). La perpendiculaire à $(d)$ passant par le point est tracée en pointillés.
    \item On reporte la même distance de l'autre côté de $(d)$, \textrr{en utilisant le compas}.
    \item On nomme le point symétrique avec une apostrophe
\end{compactitem}
\item Quand on est aidé d'un quadrillage, on place les points précisément sur les coins des carreaux. Et les traits suivent les lignes du quadrillage. Il n'y a aucune excuse pour de l'approximation.
\end{compactitem}

\end{Indication}

\end{document}

